\documentclass[11pt]{article}

\usepackage[T1]{fontenc}
\usepackage{lmodern}
\usepackage{amsmath,amssymb,amsthm,mathtools}
\usepackage{geometry}
\usepackage{hyperref}
\usepackage{enumitem}
\usepackage{booktabs}
\usepackage{microtype}

\geometry{margin=1in}
\hypersetup{
  hidelinks,
  pdftitle={The Modal Triplet Theory Program D2: Matter--Antimatter Orientation and the Baryogenesis Source Contract},
  pdfauthor={Peter Nero}
}

\newtheorem{theorem}{Theorem}
\theoremstyle{definition}
\newtheorem{definition}{Definition}
\theoremstyle{remark}
\newtheorem{remark}{Remark}

\title{The Modal Triplet Theory Program D2:\\
Matter--Antimatter Orientation and the Baryogenesis Source Contract\\
\large A Conditional Interpretation of the Finite MTT Carrier}
\author{Peter Nero}
\date{July 2026\\Version 2}

\begin{document}
\maketitle

\begin{abstract}
This paper re-evaluates the proposal that the cosmic matter--antimatter
asymmetry is a discrete orientation selection in Modal Triplet Theory (MTT).
The current finite MTT construction supplies an explicit
particle--antiparticle carrier and exact representation-theoretic consistency
checks.  That is substantial structural progress, but it is not a source
theorem for a nonzero baryon abundance.

We prove a finite carrier non-selection theorem: the same conjugation-symmetric
particle--antiparticle space and the same baryon-sign observable admit a
charge-conjugation-symmetric state with zero asymmetry and branch states with
either sign of asymmetry.  Carrier geometry, anomaly cancellation, and an
available CP-odd phase therefore do not select the state or its population
weights.  We distinguish two lawful continuations.  A branch-state route may
postulate or derive a selected asymmetric upper state and then prove its
descent.  A dynamical route must derive baryon- or lepton-number-changing
interactions, CP-odd source terms, nonequilibrium transport or an explicitly
declared alternative, and washout control from the selected action.

Discrete orientation remains a useful organizing hypothesis: it can label
the conjugate sectors and may participate in a future source mechanism.
However, the present result is a carrier theorem plus a completion and
falsification contract.  It does not derive the observed baryon-to-entropy
ratio, eliminate the need for baryogenesis dynamics, predict the absence of
antimatter domains, or turn anomaly cancellation into a population-selection
law.
\end{abstract}

\section*{Revision note for this edition}
\addcontentsline{toc}{section}{Revision note for this edition}
\begin{description}[leftmargin=2.6cm,style=nextline]
\item[Supersedes.]
Version~1.0, Zenodo record 18380140.

\item[Reason.]
The first edition treated orientation-related survivors as if their unequal
global abundance followed from refinement stability, anomaly cancellation,
and irreversibility.  The current corpus establishes a finite
particle--antiparticle carrier, but it does not establish that source map.

\item[Resolution.]
Version~2 separates carrier, state, dynamics, and observable.  It proves the
carrier non-selection result, reclassifies orientation selection as a
conditional hypothesis, restores the standard baryogenesis obligations, and
states exact promotion and falsification gates.

\item[Retained content.]
Matter and antimatter may be represented by conjugate orientation sectors of
one finite carrier, and a selected orientation variable may be useful in a
future physical mechanism.

\item[Open boundary.]
No current MTT theorem selects a nonzero baryon-number expectation, an
asymmetric branch probability, a baryogenesis source and washout history, or
the observed numerical baryon-to-entropy ratio.
\end{description}

\section*{How to read this paper}
\addcontentsline{toc}{section}{How to read this paper}
The paper is organized around one distinction:
\begin{quote}
\emph{having matter and antimatter slots is not the same as selecting how much
of each slot the universe occupies.}
\end{quote}
Section~1 places the question at the correct claim tier.  Section~2 explains
the exact finite carrier and the meaning of orientation.  Section~3 proves
that the carrier does not select an asymmetric state.  Sections~4--6 separate
anomaly constraints, CP data, and irreversibility from an actual abundance
source.  Sections~7 and 8 formulate two possible completion routes and their
required calculations.  The final sections compare the proposal with
baryogenesis, state observational gates, and record what would falsify the
MTT interpretation.

\paragraph{The central picture in plain language.}
A theatre can contain two labelled sections without deciding where the
audience sits.  Likewise, a finite Hilbert space can contain particle and
antiparticle sectors without choosing a density matrix, a cosmological state,
or a history measure.  Orientation labels the sections.  A source theorem
must still determine the occupancy and its evolution.

\tableofcontents
\newpage

\section{Question, Scope, and Current Tier}

The observed universe contains a small but nonzero excess of baryons over
antibaryons.  In conventional cosmology this is summarized by a normalized
quantity such as
\[
 \eta_B=\frac{n_B-n_{\bar B}}{s},
\]
where $s$ is the entropy density.  Explaining the sign and magnitude of
$\eta_B$ requires more than naming particle and antiparticle states.  One
must specify the state, the number observable, the dynamics that changes or
preserves that number, and the cosmological history in which the expectation
is evaluated.

Program D2 asks whether the MTT upper geometry can source those data.  The
answer is currently tiered:
\begin{enumerate}[label=\textbf{D2.\arabic*.}]
\item An explicit finite particle--antiparticle carrier exists.
\item Conjugation and orientation data can be represented on that carrier.
\item No selected state or population law producing $\eta_B\neq0$ has been
derived from that structure.
\item No complete same-source baryogenesis or asymmetric-branch calculation
has been executed.
\end{enumerate}

This paper therefore does not claim that MTT has solved baryogenesis.  It
identifies the exact structural foothold, proves the immediate non-selection
boundary, and turns the remaining idea into a testable research contract.

\section{The Finite Carrier and the Meaning of Orientation}

\subsection{What the current MTT construction supplies}

The current finite-geometry authority supplies a 96-dimensional
three-family particle--antiparticle bimodule for
\[
 \mathcal A_F=\mathbb C\oplus\mathbb H\oplus M_3(\mathbb C).
\]
Its matrix execution verifies the declared algebra action, star structure,
real and grading signs, order-zero relation, and a structural order-one
condition on the selected channel incidence.  This gives a concrete home for
particle and antiparticle slots.  The unit incidence coefficients are
structural witnesses, not measured Yukawa values or cosmic population
weights.

\subsection{Orientation as a typed label}

\begin{definition}[Conjugate orientation sectors]
Let
\[
 \mathcal H_F=\mathcal H_{+}\oplus\mathcal H_{-}
\]
and let $J_F$ be an antiunitary real-structure map exchanging the two sectors.
The signs $+$ and $-$ are called conjugate orientation labels when they
distinguish the two sectors and are reversed by $J_F$.
\end{definition}

This definition is intentionally modest.  Depending on the realization, the
label may encode particle versus antiparticle, a conjugate representation, a
sheet orientation, or another discrete involution.  Equality of labels is
not equality of particles, and the existence of the involution is not yet a
cosmological selection law.

\subsection{The baryon-sign observable}

To ask an abundance question, introduce a self-adjoint observable
\[
 B_F=
 \begin{pmatrix}
 I_{+} & 0\\
 0 & -I_{-}
 \end{pmatrix},
\qquad
 J_F B_F J_F^{-1}=-B_F.
\]
The expectation in a normalized state $\rho$ is
\[
 \langle B_F\rangle_\rho=\operatorname{Tr}(\rho B_F).
\]
The carrier and $B_F$ make the question well typed.  The density operator
$\rho$ is still independent data.

\section{The Carrier Non-Selection Theorem}

\begin{theorem}[A conjugate carrier does not select an asymmetry]
\label{thm:carrier-nonselection}
Let $\mathcal H_F=\mathcal H_+\oplus\mathcal H_-$ have nonzero finite
dimensions, let $J_F$ exchange the two sectors, and let $B_F$ have eigenvalue
$+1$ on $\mathcal H_+$ and $-1$ on $\mathcal H_-$.  Then the same fixed data
$(\mathcal H_F,J_F,B_F)$ admit normalized states $\rho_0,\rho_+,\rho_-$ with
\[
 \operatorname{Tr}(\rho_0 B_F)=0,\qquad
 \operatorname{Tr}(\rho_+ B_F)=1,\qquad
 \operatorname{Tr}(\rho_- B_F)=-1.
\]
Consequently, carrier geometry and conjugation data alone determine neither
the sign nor the magnitude of matter--antimatter asymmetry.
\end{theorem}

\begin{proof}
Choose unit vectors $u_\pm\in\mathcal H_\pm$ and define
\[
 \rho_\pm=|u_\pm\rangle\langle u_\pm|.
\]
Their expectations are respectively $+1$ and $-1$.  The state
\[
 \rho_0=\frac12
 \left(
 |u_+\rangle\langle u_+|
 +
 |J_Fu_+\rangle\langle J_Fu_+|
 \right)
\]
is normalized and has equal weight in conjugate sectors, hence zero
$B_F$ expectation.  All three states use the same carrier, involution, and
observable.
\end{proof}

\paragraph{What the theorem means.}
The result does not say that orientation is irrelevant.  It says that
orientation belongs to the kinematic inventory, while asymmetry belongs to a
state or history.  A derivation must therefore explain why one state is
selected, how its expectation is transported to four-dimensional particle
number, and why later dynamics does not erase it.

\subsection{A two-state worked example}

Take $\mathcal H_F=\mathbb C^2$, $B_F=\sigma_z$, and let conjugation exchange
the standard basis vectors.  The maximally mixed state gives
$\langle B_F\rangle=0$.  A pure upper-component state gives $+1$.  A diagonal
state
\[
 \rho_p=
 \begin{pmatrix}
 p&0\\0&1-p
 \end{pmatrix}
\]
gives $2p-1$, so every value in $[-1,1]$ occurs without changing any carrier
geometry.  This is why an orientation label cannot be used as a hidden
continuous fit for the observed abundance: the missing object is the source
of $p$.

\section{Why Anomaly Cancellation Is Not Population Selection}

Gauge and mixed anomalies are consistency conditions on chiral
representations and their quantum currents.  They determine whether a gauge
theory can be consistently quantized; they do not count how many particles
occupy a representation in a cosmological state.  The anomaly polynomial is
fixed by representation data before a density matrix such as $\rho_p$ is
chosen.

This distinction blocks a central shortcut in version~1.  If an orientation
assignment made the quantum gauge theory anomalous, that assignment would be
inadmissible as a theory.  But once an anomaly-free representation content is
fixed, anomaly cancellation does not choose a positive rather than negative
expectation of $B_F$.  In the Standard Model, the relevant anomaly
cancellations occur across the chiral multiplets of a generation; they are
not a cosmological mechanism that deletes antiparticle states.

\begin{remark}[Correct use of anomaly data]
Anomaly freedom is a gate on a proposed source theory.  It may rule out a
candidate interaction or representation assignment.  Passing the gate leaves
the state-selection and transport problem intact.
\end{remark}

\section{What a CP Phase Can and Cannot Do}

A CP-odd phase can enter interference terms and permit different rates for
conjugate processes.  That makes it relevant to baryogenesis, but not
sufficient by itself.  The classic Sakharov analysis identifies
baryon-number violation, C and CP violation, and departure from thermal
equilibrium as the standard conditions for dynamically generating an
asymmetry from a symmetric initial state \cite{Sakharov1967}.  Electroweak
calculations further show that the observed CKM source is not by itself large
enough to account for the cosmic asymmetry in the minimal Standard Model
\cite{HuetSather1995}.

The MTT corpus contains finite-phase and orientation clues associated with the
selected q79 branch.  At the current authority tier, those clues do not
establish a physical map to all CKM angles, a baryon-number source, or
$\eta_B$.  A shared-circle holonomy becomes relevant only after one proves an
intertwiner
\[
 L_{\rm shared}
 \longrightarrow
 \text{physical CP-odd source on }\mathcal H_F
\]
that preserves the selected connection and action and is evaluated in the
same cosmological state.  Merely observing that both structures contain a
phase is not such a theorem.

\section{Irreversibility, CPT, and Branch Selection}

\subsection{Persistence is not production}

Suppose a branch already has $\langle B_F\rangle\neq0$ and the effective
dynamics approximately conserves that expectation after some epoch.
Irreversibility or washout suppression can then explain persistence.  It
cannot explain the original sign or magnitude unless the same mechanism also
emits the asymmetric state or source term.  Production and persistence are
different proof obligations.

\subsection{CPT-compatible asymmetric states}

An asymmetric cosmological state need not mean that the local action violates
CPT.  A symmetry of the dynamics need not be a symmetry of a particular
state or boundary condition.  Conversely, one cannot claim that a vague
failure of ``global admissibility'' selects matter while local CPT remains
intact.  The claim requires a constructed local quantum field theory, its CPT
map, the selected state, and a proof of how the state descends.

The safe current statement is:
\begin{quote}
MTT may eventually select a CPT-compatible asymmetric branch state, but no
current theorem establishes that selection or its probability.
\end{quote}

\subsection{Two branches are not automatically two universes}

The existence of conjugate mathematical states says only that both are
allowed by the carrier.  Interpreting them as two physical universes, as two
time orientations, or as one realized and one hidden branch requires an
additional ontology and selection law.  Program D2 does not assume any of
those readings.

\section{Completion Route A: A Selected Asymmetric Branch State}

The first lawful route is kinematic.  One selected upper state
$\Omega_{\rm up}$ may descend to a finite state $\rho_F$ and then to an
effective four-dimensional state.  To close this route, one must provide:
\begin{enumerate}[label=\textbf{A\arabic*.}]
\item a normalized upper state or history measure;
\item a selected branch rule that does not use the observed asymmetry as its
selector;
\item an intertwiner from the upper orientation observable to $B_F$;
\item a state-preserving descent to the four-dimensional baryon current;
\item a proof that subsequent evolution preserves the sign and yields a
controlled magnitude;
\item a probability or typicality statement if the theory claims that the
observed branch is preferred rather than merely allowed.
\end{enumerate}

This route can consistently say that the asymmetry is an initial or global
state property.  It cannot call that property a derived prediction until the
source and branch rule are supplied.

\section{Completion Route B: Selected Baryogenesis Dynamics}

The second route derives an asymmetry dynamically from a declared earlier
state.  At effective level, many mechanisms can be organized schematically by
a transport equation
\[
 \dot n_B+3Hn_B=S_B[\Phi_{\rm CP},\Gamma_{\Delta B},T,\ldots]
 -\Gamma_{\rm wash}\,n_B .
\]
Here $S_B$ is a CP-odd baryon source, $\Gamma_{\Delta B}$ represents
baryon- or lepton-number-changing processes, and $\Gamma_{\rm wash}$ removes
the generated asymmetry.  This equation is a diagnostic template, not an MTT
derivation.

To promote it, MTT must derive from one selected action:
\begin{enumerate}[label=\textbf{B\arabic*.}]
\item the relevant number current and its violation law;
\item the CP-odd invariant and its coupling to the reactions;
\item the nonequilibrium background, or a precisely stated alternative to
the usual Sakharov setting;
\item finite-temperature or nonequilibrium rates with uncertainties;
\item washout and spectator processes;
\item the transition from any lepton asymmetry to baryon asymmetry when that
route is used;
\item the final $\eta_B$ prediction without choosing inputs from the target
value.
\end{enumerate}

The current open upper-action, nonperturbative QFT, and no-knob Standard Model
value blockers sit upstream of this computation.  Closing the finite carrier
does not bypass them.

\section{Relation to the MTT Program}

\subsection{Program B2 and B3}

Program B2 supplies a conditional language for gauge structure; Program B3
supplies a conditional language for discrete survivor structure.  D2 may use
their objects once their hypotheses are checked, but it cannot infer a
population law from the words ``gauge'' or ``discrete.''  Representation,
state, and abundance remain different types.

\subsection{Program D1}

Program D1 distinguishes an availability mask from a covariant source.
Program D2 has the same logical shape.  The particle--antiparticle carrier is
the mask of available sectors; the selected state or baryogenesis action is
the source.  In both papers, the missing physical map cannot be supplied by
renaming the structural data.

\subsection{The finite Standard Model carrier}

The explicit 96-dimensional carrier is the principal new foothold unavailable
to version~1.  It makes the orientation hypothesis calculable: a future
operator can now be written on named particle and antiparticle slots.  Its
current achievement is exact representation compatibility.  Physical finite
Dirac entries, full source-owned flavor data, nonperturbative QFT transport,
and cosmological state selection remain separate.

\section{Observational and Numerical Contract}

A completed Program D2 should report at least:
\begin{enumerate}[label=\textbf{O\arabic*.}]
\item the convention for $B$, $L$, and $B-L$;
\item the selected initial state and cosmological background;
\item the complete source and washout equations;
\item the predicted sign and interval for $\eta_B$;
\item the relation to light-element abundances and cosmic microwave
background constraints;
\item any predicted antimatter-domain or isocurvature signature;
\item an uncertainty budget and a held-out comparison protocol.
\end{enumerate}

The numerical value of $\eta_B$ may be used for final comparison, not as the
coefficient that chooses an orientation, phase, branch, or rate.  If one to
three physical primitives are retained, they must be declared before the
comparison and shown not to be disguised replay coordinates.

\section{Falsification and Non-Promotion Rules}

The orientation program fails at its claimed tier if any of the following
occurs:
\begin{enumerate}[label=\textbf{F\arabic*.}]
\item the proposed orientation operator does not intertwine with the physical
particle--antiparticle carrier;
\item the selected upper state descends to zero or the wrong-sign asymmetry;
\item the source action violates gauge consistency, locality, unitarity, or
the declared CPT contract;
\item washout removes the generated signal;
\item the calculated asymmetry depends on choosing a branch or coefficient
from the observed $\eta_B$;
\item a claimed prediction lacks an uncertainty-controlled comparison with
cosmological data.
\end{enumerate}

The following statements are not promotion evidence:
\begin{itemize}
\item the carrier contains antiparticle slots;
\item a finite phase is numerically near a measured CP coordinate;
\item an anomaly-free representation exists;
\item one asymmetric state can be written down;
\item irreversible evolution preserves an asymmetry inserted by hand.
\end{itemize}

\section{Conclusion}

Program D2 now has a firmer and narrower foundation.  MTT possesses a concrete
finite particle--antiparticle carrier on which orientation, conjugation, and
number observables can be represented.  The carrier non-selection theorem
shows exactly why that achievement does not yet explain the cosmic abundance:
the same finite data support symmetric and oppositely asymmetric states.

Discrete orientation therefore survives as a candidate source coordinate,
not as a completed explanation.  The next scientific advance must select
either an asymmetric upper state with a lawful descent or a baryogenesis
dynamics with number violation, CP-odd sourcing, nonequilibrium transport,
and washout control.  Both routes can be tested on the existing finite
carrier.  Neither is supplied by anomaly cancellation, a phase label, or
irreversibility alone.

This correction is constructive rather than destructive.  It replaces an
unverifiable narrative with a finite operator problem and a clear exit
certificate.  A successful future version will not merely say that matter is
the selected orientation; it will calculate why the selected state has the
observed sign and magnitude and show that the result survives the complete
four-dimensional dynamics.

\section*{Computational Evidence and Reproducibility}
\addcontentsline{toc}{section}{Computational Evidence and Reproducibility}
The finite carrier result used here is archived in the curated MTT results
repository under authority A48:
\begin{quote}\small\raggedright
\url{https://github.com/PeterNero/mtt-results-repro/tree/main/release/authority/A48}.
\end{quote}
Authority A48 establishes the finite representation carrier and explicitly
does not promote structural incidence coefficients to physical Yukawa or
population values.  The no-knob Standard Model value and precision-equivalence
frontier remains open under blocker B.SM.02 in the MTT research ledger.

\begin{thebibliography}{99}

\bibitem{Sakharov1967}
A.~D.~Sakharov,
``Violation of CP invariance, C asymmetry, and baryon asymmetry of the
universe,''
\emph{JETP Letters} \textbf{5}, 24--27 (1967);
\href{https://doi.org/10.1070/PU1991v034n05ABEH002497}
{doi:10.1070/PU1991v034n05ABEH002497}.

\bibitem{KobayashiMaskawa1973}
M.~Kobayashi and T.~Maskawa,
``CP-violation in the renormalizable theory of weak interaction,''
\emph{Progress of Theoretical Physics} \textbf{49}, 652--657 (1973),
\href{https://doi.org/10.1143/PTP.49.652}
{doi:10.1143/PTP.49.652}.

\bibitem{tHooft1976}
G.~'t Hooft,
``Symmetry breaking through Bell--Jackiw anomalies,''
\emph{Physical Review Letters} \textbf{37}, 8--11 (1976),
\href{https://doi.org/10.1103/PhysRevLett.37.8}
{doi:10.1103/PhysRevLett.37.8}.

\bibitem{KuzminRubakovShaposhnikov1985}
V.~A.~Kuzmin, V.~A.~Rubakov, and M.~E.~Shaposhnikov,
``On the anomalous electroweak baryon-number non-conservation in the early
universe,''
\emph{Physics Letters B} \textbf{155}, 36--42 (1985),
\href{https://doi.org/10.1016/0370-2693(85)91028-7}
{doi:10.1016/0370-2693(85)91028-7}.

\bibitem{HuetSather1995}
P.~Huet and E.~Sather,
``Electroweak baryogenesis and Standard Model CP violation,''
\emph{Physical Review D} \textbf{51}, 379--394 (1995),
\href{https://doi.org/10.1103/PhysRevD.51.379}
{doi:10.1103/PhysRevD.51.379}.

\bibitem{Greenberg2002}
O.~W.~Greenberg,
``CPT violation implies violation of Lorentz invariance,''
\emph{Physical Review Letters} \textbf{89}, 231602 (2002),
\href{https://doi.org/10.1103/PhysRevLett.89.231602}
{doi:10.1103/PhysRevLett.89.231602}.

\end{thebibliography}

\end{document}
